%! Matrices as Functions — Unit 4, Lesson 4.13 — Student Packet Cover Sheet
\documentclass[10pt]{article}
\usepackage{precalculus-article}
\usepackage{precalculus-boxes}
\usepackage{ltablex}
\keepXColumns

\begin{document}

% Full-bleed navy banner
\begin{tikzpicture}[remember picture, overlay]
  \fill[navy] (current page.north west)
    rectangle ([yshift=-0.9in]current page.north east);
\end{tikzpicture}
\vspace{-0.8in}

\begin{center}
  {\color{white}\textbf{\LARGE AP Precalculus}}\\[4pt]
  {\color{skymid}\large\textbf{Unit 4: Functions Involving Parameters, Vectors, and Matrices}}\\[6pt]
  {\color{white}\textbf{\large Lesson 4.13 \quad Matrices as Functions}}
\end{center}

\vspace{0.22in}
\namedateperiod
\vspace{0.18in}

\begin{learningtargetbox}
\small
\begin{enumerate}[leftmargin=1.5em, itemsep=5pt]
  \item I can find the matrix associated with a linear transformation using the images of the unit vectors.
  \item I can apply the rotation matrix $R(\theta)$ to rotate vectors counterclockwise by angle $\theta$.
  \item I can compose two linear transformations by multiplying their matrices in the correct order.
  \item I can find the inverse of a linear transformation using the matrix inverse.
\end{enumerate}
\end{learningtargetbox}

\vspace{0.14in}

\begin{tocbox}
\small
\renewcommand{\arraystretch}{1.5}
\begin{tabularx}{\linewidth}{@{} c l X r @{}}
\rowcolor{sky}
\textbf{\#} & \textbf{Component} & \textbf{Description} & \textbf{Score} \\
\midrule
1 & Warm-Up        & Matrix--vector products and the $2\times2$ inverse & \blank{1.2cm} \\
2 & Guided Notes   & Reading the matrix; rotation; composition; inverse  & \blank{1.2cm} \\
3 & Group Activity & Tiered transformation problems (Tiers R / A / E)    & \blank{1.2cm} \\
4 & Exit Ticket    & Three short transformation questions                 & \blank{1.2cm} \\
5 & Homework       & Six practice problems + AP-style MC                  & \blank{1.2cm} \\
\midrule
  & \multicolumn{2}{r}{\textbf{Total}} & \blank{1.2cm} \\
\end{tabularx}
\end{tocbox}

\end{document}
